\section{Related Work}
\label{sec:related_work}

\subsection{Remaining Useful Life Point Prognostics}
Prognostic health management of aero-engines was popularized by the NASA C-MAPSS dataset release~\cite{saxena2008damage}, which formalized run-to-failure simulation across varied operational flight conditions and fault modes. Standard point predictors model RUL using deep neural networks or gradient-boosted decision trees. While the C-MAPSS asymmetric scoring metric heavily penalizes overestimation ($d > 0$), base models are routinely optimized under symmetric loss objectives. Furthermore, point predictions provide no confidence intervals, depriving operators of actionable risk metrics.

\subsection{Conformal Uncertainty Quantification in PHM}
Conformal prediction~\cite{javanmardi2023conformal} has emerged as an authoritative distribution-free uncertainty framework. Conformalized Quantile Regression (CQR) was adapted to aero-engines by Diao et al.~\cite{diao2026turbofan} using LSTM networks on C-MAPSS FD001 and FD002, and by Robinson~\cite{robinson2026riskaware} across all four C-MAPSS subsets. However, standard split conformal methods assume observation exchangeability. Because sensor telemetry is autocorrelated within an engine, pooling raw cycles violates exchangeability and over-represents long-lived engines. Lee, Barber, and Willett~\cite{lee2023hierarchical} established hierarchical conformal inference for clustered data, placing $\frac{1}{m+1}$ probability mass at $+\infty$ to target valid coverage for unseen clusters under exchangeable group sampling.

IntelliTwin uses an engine-weighted conformal construction motivated by hierarchical conformal prediction. Formal coverage results require the hierarchical sampling and exchangeability assumptions stated in the underlying theory. Equal engine weighting does not establish these assumptions for ordered, lifetime-stratified degradation trajectories. This study therefore evaluates marginal and subgroup coverage empirically. It does not claim guaranteed coverage for each engine, lifecycle stage, operating regime, complete trajectory, or adaptively selected maintenance time. Nominal 90% coverage is a target, not a guaranteed realized coverage proportion. Maintenance outcomes are evaluated separately under the frozen simulator.

\subsection{Conditional Reliability Diagnostics and Maintenance Economics}
Recent studies emphasize that marginal coverage obscures severe conditional vulnerabilities. Yang et al.~\cite{yang2026empirical} audited empirical calibration across operating-regime shifts in rotating bearings, revealing acute localized undercoverage. BenAbdennour~\cite{benabdennour2026grouped} demonstrated grouped validation and interval-safety calibration across degradation domains, showing non-linear expansion of uncertainty during terminal wear stages. 

Concurrently, uncertainty-aware maintenance scheduling has been investigated via simulation and reinforcement learning~\cite{chen2022data,zhu2025predictive,walia2026uncertainty}. Cortes-Gomez et al.~\cite{cortesgomez2025utility} introduced Utility-Directed Conformal Prediction (UDCP), demonstrating that tailoring nonconformity scores to downstream decision utility substantially reduces operational loss. IntelliTwin synthesizes these domains, evaluating whether asymmetric decision-aware score inversion can improve sequential maintenance utility while preserving subgroup reliability and coverage safety under the frozen evaluation protocol.
